\documentclass[a4paper,12pt]{article}

\usepackage[T2A]{fontenc}
\usepackage[utf8]{inputenc}
\usepackage[russian]{babel}
\usepackage{amsmath,amssymb,amsthm,mathtools}
\usepackage{helvet}
\renewcommand{\familydefault}{\sfdefault}
\usepackage{icomma}
\usepackage[
  a4paper,
  left=18mm,
  right=18mm,
  top=18mm,
  bottom=20mm,
  headheight=14pt
]{geometry}
\usepackage{microtype}
\usepackage{tikz}
\usetikzlibrary{calc}
\usetikzlibrary{arrows.meta,positioning,shapes.geometric}
\usepackage{tkz-euclide}
\usepackage{pgfplots}
\pgfplotsset{compat=1.18}
\usepackage{tabularx,booktabs,longtable}
\usepackage{enumitem}
\usepackage[hidelinks]{hyperref}
\usepackage{parskip}
\usepackage{fancyhdr}
\usepackage{setspace}
\usepackage{xcolor}

\definecolor{accent}{HTML}{008080}
\definecolor{darkaccent}{HTML}{004D4D}
\definecolor{textgray}{HTML}{333333}
\definecolor{softgray}{HTML}{727B7B}

\color{textgray}
\onehalfspacing
\setlength{\parindent}{0pt}

\pagestyle{fancy}
\fancyhf{}
\fancyhead[L]{\small Сечения куба}
\fancyhead[R]{\small Ксения Клыкова, 10 класс}
\fancyfoot[C]{\small\thepage}
\renewcommand{\headrulewidth}{0.3pt}

\setlist[itemize]{leftmargin=6mm,itemsep=2pt,topsep=3pt}
\setlist[enumerate]{leftmargin=7mm,itemsep=2pt,topsep=3pt}

\newcommand{\term}[1]{\textbf{\color{darkaccent}#1}}
\newcommand{\Plane}{\ensuremath{\pi}}

\tikzset{
  flowstart/.style={
    ellipse,
    draw=darkaccent,
    line width=0.9pt,
    align=center,
    minimum width=48mm,
    minimum height=10mm,
    inner sep=2.5mm,
    font=\small
  },
  flowaction/.style={
    rounded corners=2mm,
    rectangle,
    draw=accent,
    line width=0.9pt,
    align=center,
    text width=88mm,
    minimum height=11mm,
    inner sep=2.5mm,
    font=\small
  },
  flowdecision/.style={
    diamond,
    aspect=2.4,
    draw=darkaccent,
    line width=0.9pt,
    align=center,
    text width=46mm,
    inner sep=1.5mm,
    font=\small
  },
  flowarrow/.style={
    -{Stealth[length=2.8mm,width=2mm]},
    line width=0.9pt,
    color=textgray
  }
}

\begin{document}

% Источник: доска занятия со схемами нескольких построений сечений.
% :contentReference[oaicite:0]{index=0}

% ============================================================
% ТИТУЛЬНЫЙ ЛИСТ
% ============================================================

\begin{titlepage}
\thispagestyle{empty}
\centering

\vspace*{22mm}

{\Large\color{darkaccent}\textbf{Чек-лист}}\par
\vspace{5mm}
{\LARGE\textbf{Сечения куба}}\par
\vspace{2mm}
{\large Построение сечений по трём точкам}\par

\vspace{14mm}

{\large Ксения Клыкова, 10 класс}\par
\vspace{5mm}
{\normalsize Дата: \today}\par

\vfill

{\large\textbf{Лёвин Артём Александрович}}\par
\vspace{2mm}
{\normalsize эксклюзивно для Ксении Клыковой}\par

\vspace{24mm}

\begin{tikzpicture}[remember picture,overlay]
\draw[
  accent!70,
  line width=1.1pt
]
  ($(current page.south west)+(16mm,18mm)$)
  ..
  controls
  ($(current page.south west)+(72mm,30mm)$)
  and
  ($(current page.south east)+(-75mm,10mm)$)
  ..
  ($(current page.south east)+(-16mm,20mm)$);
\end{tikzpicture}

\end{titlepage}

% ============================================================
% КРАТКИЙ ЧЕК-ЛИСТ ПО ТЕМЕ
% ============================================================

\section*{1. Главное по теме}

\term{Сечение куба} --- это многоугольник, полученный при пересечении куба
с секущей плоскостью.

\term{Три неколлинеарные точки} однозначно задают секущую плоскость \Plane.

\term{След секущей плоскости на грани} --- прямая, по которой плоскость \Plane
пересекает плоскость этой грани.

\term{Базовые идеи построения}

\begin{itemize}
\item две точки можно соединять, если они лежат в одной грани;
\item если две точки принадлежат секущей плоскости, то вся прямая через них
тоже принадлежит этой плоскости;
\item общее ребро двух соседних граней служит переходом от одной грани к другой;
\item внешняя вспомогательная точка допустима, если она получена как пересечение
двух прямых, лежащих в одной плоскости;
\item цель построения --- найти вершины сечения на рёбрах куба и замкнуть
многоугольник.
\end{itemize}

\term{Мини-памятка проверки}

\begin{itemize}
\item каждая сторона сечения лежит в одной из граней куба;
\item каждая вершина сечения лежит на ребре куба;
\item порядок вершин образует замкнутый многоугольник;
\item вспомогательные построения не дают ложных пересечений.
\end{itemize}

% ============================================================
% АЛГОРИТМ
% ============================================================

\clearpage
\thispagestyle{plain}

\begin{center}
{\Large\color{darkaccent}\textbf{Алгоритм построения сечения}}
\end{center}

\vspace{4mm}

\begin{center}
\begin{tikzpicture}[node distance=7mm and 16mm]
\node[flowstart] (start)
{Даны три точки\\$P,Q,R$};

\node[flowaction,below=of start] (step1)
{Определите, какие две из них лежат в одной грани,\\и проведите первый след плоскости};

\node[flowdecision,below=9mm of step1] (dec1)
{Сечение уже замкнулось?};

\node[flowaction,below=9mm of dec1] (step2)
{Выберите соседнюю грань\\и найдите общее ребро двух граней};

\node[flowaction,below=of step2] (step3)
{Продлите известный след до пересечения\\с общим ребром. Получите вспомогательную точку};

\node[flowaction,below=of step3] (step4)
{Соедините вспомогательную точку\\с известной точкой секущей плоскости на новой грани};

\node[flowaction,below=of step4] (step5)
{Найдите новую вершину сечения на ребре куба\\и продолжите построение};

\node[flowstart,below=of step5] (finish)
{Многоугольник замкнут:\\сечение построено};

\draw[flowarrow] (start) -- (step1);
\draw[flowarrow] (step1) -- (dec1);
\draw[flowarrow] (dec1) -- node[right,font=\small]{нет} (step2);
\draw[flowarrow] (step2) -- (step3);
\draw[flowarrow] (step3) -- (step4);
\draw[flowarrow] (step4) -- (step5);
\draw[flowarrow] (step5.west) -- ++(-24mm,0) |- (dec1.west);
\draw[flowarrow] (dec1.east) -- ++(24mm,0) |- node[pos=0.2,above,font=\small]{да} (finish.east);
\end{tikzpicture}
\end{center}

% ============================================================
% РАЗОБРАННЫЕ ЗАДАНИЯ
% ============================================================

\clearpage
\section*{2. Разобранные задания с чертежами}

Во всех рисунках используется стандартная разметка куба
$ABCD\,EFGH$:

\[
ABCD \text{ --- передняя грань,}
\qquad
EFGH \text{ --- задняя грань.}
\]

Ниже приведены все восстановленные по содержанию занятия типовые
разобранные построения. Для каждого задания дан чертёж и короткая
проверка корректности.

% ------------------------------------------------------------
% ЗАДАНИЕ 1
% ------------------------------------------------------------

\subsection*{Задание 1}

Пусть
\[
P\in HE,\qquad Q\in FG,\qquad R\in BC.
\]
Требуется построить сечение куба плоскостью $\Plane=PQR$.

\begin{minipage}[t]{0.53\linewidth}
\begin{center}
\begin{tikzpicture}[scale=1.0]
\tkzSetUpLine[line width=0.9pt,color=textgray]
\tkzSetUpPoint[color=accent,fill=accent,size=3]
\tkzSetUpLabel[color=black,font=\small]

\tkzDefPoint(0,2){A}
\tkzDefPoint(3,2){B}
\tkzDefPoint(3,0){C}
\tkzDefPoint(0,0){D}
\tkzDefPoint(1,3){E}
\tkzDefPoint(4,3){F}
\tkzDefPoint(4,1){G}
\tkzDefPoint(1,1){H}

\coordinate (P) at ($(H)!0.65!(E)$);
\coordinate (Q) at ($(F)!0.55!(G)$);
\coordinate (R) at ($(B)!0.75!(C)$);
\coordinate (S) at ($(D)!0.45!(A)$);

\tkzDrawSegments(A,B B,C C,D D,A A,E B,F C,G D,H)
\tkzDrawSegments[dashed,color=softgray](E,F F,G G,H H,E)

\draw[darkaccent,line width=1.4pt] (P)--(Q)--(R)--(S)--cycle;

\fill[accent] (P) circle (2pt);
\fill[accent] (Q) circle (2pt);
\fill[accent] (R) circle (2pt);
\fill[accent] (S) circle (2pt);

\tkzLabelPoints[above left](A)
\tkzLabelPoints[above right](B,F)
\tkzLabelPoints[below right](C,G)
\tkzLabelPoints[below left](D,H)
\tkzLabelPoints[above](E)
\node[left] at (P) {$P$};
\node[right] at (Q) {$Q$};
\node[right] at (R) {$R$};
\node[left] at (S) {$S$};
\end{tikzpicture}
\end{center}
\end{minipage}
\hfill
\begin{minipage}[t]{0.43\linewidth}
\textbf{Итог:} сечение --- четырёхугольник
\[
SPQR.
\]

\textbf{Построение.}
\begin{enumerate}
\item $Q$ и $R$ лежат в правой грани $BCGF$, поэтому проводим $QR$.
\item $P$ и $Q$ лежат в задней грани $EFGH$, поэтому проводим $PQ$.
\item Следы секущей плоскости приходят к передней и левой граням, где
получается четвёртая вершина $S$ на ребре $DA$.
\item Соединяем $S$ с $P$.
\end{enumerate}

\textbf{Проверка.}
\begin{itemize}
\item $PQ\subset(EFGH)$, $QR\subset(BCGF)$;
\item $RS\subset(ABCD)$, $SP\subset(ADHE)$;
\item получен замкнутый четырёхугольник без самопересечений.
\end{itemize}
\end{minipage}

\vspace{4mm}

% ------------------------------------------------------------
% ЗАДАНИЕ 2
% ------------------------------------------------------------

\subsection*{Задание 2}

Пусть
\[
P\in AE,\qquad Q\in BF,\qquad R\in CG.
\]
Требуется построить сечение куба плоскостью $\Plane=PQR$.

\begin{minipage}[t]{0.53\linewidth}
\begin{center}
\begin{tikzpicture}[scale=1.0]
\tkzSetUpLine[line width=0.9pt,color=textgray]
\tkzSetUpPoint[color=accent,fill=accent,size=3]
\tkzSetUpLabel[color=black,font=\small]

\tkzDefPoint(0,2){A}
\tkzDefPoint(3,2){B}
\tkzDefPoint(3,0){C}
\tkzDefPoint(0,0){D}
\tkzDefPoint(1,3){E}
\tkzDefPoint(4,3){F}
\tkzDefPoint(4,1){G}
\tkzDefPoint(1,1){H}

\coordinate (P) at ($(A)!0.55!(E)$);
\coordinate (Q) at ($(B)!0.55!(F)$);
\coordinate (R) at ($(C)!0.45!(G)$);
\coordinate (S) at ($(D)!0.45!(H)$);

\tkzDrawSegments(A,B B,C C,D D,A A,E B,F C,G D,H)
\tkzDrawSegments[dashed,color=softgray](E,F F,G G,H H,E)

\draw[darkaccent,line width=1.4pt] (P)--(Q)--(R)--(S)--cycle;

\fill[accent] (P) circle (2pt);
\fill[accent] (Q) circle (2pt);
\fill[accent] (R) circle (2pt);
\fill[accent] (S) circle (2pt);

\tkzLabelPoints[above left](A)
\tkzLabelPoints[above right](B,F)
\tkzLabelPoints[below right](C,G)
\tkzLabelPoints[below left](D,H)
\tkzLabelPoints[above](E)
\node[left] at (P) {$P$};
\node[right] at (Q) {$Q$};
\node[right] at (R) {$R$};
\node[left] at (S) {$S$};
\end{tikzpicture}
\end{center}
\end{minipage}
\hfill
\begin{minipage}[t]{0.43\linewidth}
\textbf{Итог:} сечение --- четырёхугольник
\[
PQRS.
\]

\textbf{Построение.}
\begin{enumerate}
\item $P$ и $Q$ лежат в верхней грани $ABFE$, поэтому проводим $PQ$.
\item $Q$ и $R$ лежат в правой грани $BCGF$, поэтому проводим $QR$.
\item Последняя вершина получается на ребре $DH$:
\[
S\in DH.
\]
\item Соединяем $S$ с $P$.
\end{enumerate}

\textbf{Проверка.}
\begin{itemize}
\item $PQ\subset(ABFE)$, $QR\subset(BCGF)$;
\item $RS\subset(DCGH)$, $SP\subset(ADHE)$;
\item стороны идут по соседним граням, многоугольник замкнут.
\end{itemize}
\end{minipage}

\vspace{4mm}

% ------------------------------------------------------------
% ЗАДАНИЕ 3
% ------------------------------------------------------------

\subsection*{Задание 3}

Пусть
\[
P\in DA,\qquad Q\in BF,\qquad R\in CG.
\]
Требуется построить сечение куба плоскостью $\Plane=PQR$.

\begin{minipage}[t]{0.53\linewidth}
\begin{center}
\begin{tikzpicture}[scale=1.0]
\tkzSetUpLine[line width=0.9pt,color=textgray]
\tkzSetUpPoint[color=accent,fill=accent,size=3]
\tkzSetUpLabel[color=black,font=\small]

\tkzDefPoint(0,2){A}
\tkzDefPoint(3,2){B}
\tkzDefPoint(3,0){C}
\tkzDefPoint(0,0){D}
\tkzDefPoint(1,3){E}
\tkzDefPoint(4,3){F}
\tkzDefPoint(4,1){G}
\tkzDefPoint(1,1){H}

\coordinate (P) at ($(D)!0.55!(A)$);
\coordinate (Q) at ($(A)!0.265!(B)$);
\coordinate (R) at ($(B)!0.25!(F)$);
\coordinate (S) at ($(C)!0.45!(G)$);
\coordinate (T) at ($(D)!0.11!(H)$);

\tkzDrawSegments(A,B B,C C,D D,A A,E B,F C,G D,H)
\tkzDrawSegments[dashed,color=softgray](E,F F,G G,H H,E)

\draw[darkaccent,line width=1.4pt] (P)--(Q)--(R)--(S)--(T)--cycle;

\fill[accent] (P) circle (2pt);
\fill[accent] (Q) circle (2pt);
\fill[accent] (R) circle (2pt);
\fill[accent] (S) circle (2pt);
\fill[accent] (T) circle (2pt);

\tkzLabelPoints[above left](A)
\tkzLabelPoints[above right](B,F)
\tkzLabelPoints[below right](C,G)
\tkzLabelPoints[below left](D,H)
\tkzLabelPoints[above](E)

\node[left] at (P) {$P$};
\node[above] at (Q) {$Q$};
\node[right] at (R) {$R$};
\node[right] at (S) {$S$};
\node[left] at (T) {$T$};
\end{tikzpicture}
\end{center}
\end{minipage}
\hfill
\begin{minipage}[t]{0.43\linewidth}
\textbf{Итог:} сечение --- пятиугольник
\[
PQRST.
\]

\textbf{Построение.}
\begin{enumerate}
\item В передней грани сначала связываются точки $P$ и $Q$.
\item Затем через верхнюю и правую грани строятся стороны $QR$ и $RS$.
\item След секущей плоскости переходит на нижнюю грань и даёт вершину
$T$ на ребре $DH$.
\item Последняя сторона --- $TP$.
\end{enumerate}

\textbf{Проверка.}
\begin{itemize}
\item $PQ\subset(ABCD)$;
\item $QR\subset(ABFE)$, $RS\subset(BCGF)$;
\item $ST\subset(DCGH)$, $TP\subset(ADHE)$;
\item каждая сторона принадлежит своей грани, порядок вершин корректен.
\end{itemize}
\end{minipage}

\vspace{4mm}

% ------------------------------------------------------------
% ЗАДАНИЕ 4
% ------------------------------------------------------------

\subsection*{Задание 4}

Пусть
\[
P\in AB,\qquad Q\in BC,\qquad R\in CG.
\]
Дополнительно секущая плоскость проходит через точки
\[
T\in AE,\qquad U\in HE,\qquad S\in GH,
\]
образуя полный шестиугольник сечения.

\begin{minipage}[t]{0.53\linewidth}
\begin{center}
\begin{tikzpicture}[scale=1.0]
\tkzSetUpLine[line width=0.9pt,color=textgray]
\tkzSetUpPoint[color=accent,fill=accent,size=3]
\tkzSetUpLabel[color=black,font=\small]

\tkzDefPoint(0,2){A}
\tkzDefPoint(3,2){B}
\tkzDefPoint(3,0){C}
\tkzDefPoint(0,0){D}
\tkzDefPoint(1,3){E}
\tkzDefPoint(4,3){F}
\tkzDefPoint(4,1){G}
\tkzDefPoint(1,1){H}

\coordinate (P) at ($(A)!0.30!(B)$);
\coordinate (Q) at ($(B)!0.737!(C)$);
\coordinate (R) at ($(C)!0.50!(G)$);
\coordinate (S) at ($(G)!0.25!(H)$);
\coordinate (U) at ($(H)!0.789!(E)$);
\coordinate (T) at ($(A)!0.60!(E)$);

\tkzDrawSegments(A,B B,C C,D D,A A,E B,F C,G D,H)
\tkzDrawSegments[dashed,color=softgray](E,F F,G G,H H,E)

\draw[darkaccent,line width=1.4pt] (P)--(Q)--(R)--(S)--(U)--(T)--cycle;

\fill[accent] (P) circle (2pt);
\fill[accent] (Q) circle (2pt);
\fill[accent] (R) circle (2pt);
\fill[accent] (S) circle (2pt);
\fill[accent] (U) circle (2pt);
\fill[accent] (T) circle (2pt);

\tkzLabelPoints[above left](A)
\tkzLabelPoints[above right](B,F)
\tkzLabelPoints[below right](C,G)
\tkzLabelPoints[below left](D,H)
\tkzLabelPoints[above](E)

\node[above] at (P) {$P$};
\node[right] at (Q) {$Q$};
\node[right] at (R) {$R$};
\node[below] at (S) {$S$};
\node[left] at (U) {$U$};
\node[left] at (T) {$T$};
\end{tikzpicture}
\end{center}
\end{minipage}
\hfill
\begin{minipage}[t]{0.43\linewidth}
\textbf{Итог:} сечение --- шестиугольник
\[
PQRSTU.
\]

\textbf{Построение.}
\begin{enumerate}
\item Начинаем с передней и правой граней:
\[
PQ,\quad QR.
\]
\item Затем след секущей плоскости переходит на нижнюю грань:
\[
RS.
\]
\item После этого строятся задняя и левая части сечения:
\[
SU,\quad UT.
\]
\item Многоугольник замыкается стороной $TP$.
\end{enumerate}

\textbf{Проверка.}
\begin{itemize}
\item $PQ\subset(ABCD)$, $QR\subset(BCGF)$;
\item $RS\subset(DCGH)$, $SU\subset(EFGH)$;
\item $UT\subset(ADHE)$, $TP\subset(ABFE)$;
\item шесть вершин лежат на рёбрах куба, самопересечений нет.
\end{itemize}
\end{minipage}

% ============================================================
% ТИПИЧНЫЕ ОШИБКИ
% ============================================================

\clearpage
\section*{3. Частые ошибки}

\begin{itemize}
\item соединять точки без проверки, лежат ли они в одной грани;
\item продлевать прямую в случайную сторону;
\item путать вспомогательную внешнюю точку и вершину самого сечения;
\item не проверять, на каком общем ребре соседних граней нужно делать переход;
\item забывать замкнуть многоугольник;
\item не проверять принадлежность последней стороны нужной грани.
\end{itemize}

\textbf{Быстрая проверка после построения}

\begin{enumerate}
\item Назовите все вершины сечения по порядку.
\item Для каждой стороны укажите грань, в которой она лежит.
\item Для каждой вершины укажите ребро куба, на котором она расположена.
\item Убедитесь, что многоугольник замкнут и не имеет ложных пересечений.
\end{enumerate}


\end{document}